\documentclass[12pt]{article}
\usepackage[utf8]{inputenc}
\usepackage[T1]{fontenc}
\usepackage{hyperref}
\title{Bibbulman Language Resources for K-12 Curriculum Development}
\author{Compiled by Claude Sonnet 5.5 (claude-sonnet-5-5) for Andr\'as Kornai}
\date{}
\begin{document}
\maketitle

\section{Introduction}
Bibbulman (also Pipelman) is a variety of the Noongar language of south-western Western Australia. It has its own ISO code (xbp) and is listed as extinct. Its country is placed between the Pemberton area and Broke Inlet. Noongar as a whole is living: the 2016 census recorded 443 people speaking it at home, and one Wikipedia version gives under 240 speakers (2014). I found no Bibbulman-specific dictionary, course or school programme in search results. Any curriculum would draw on wider Noongar resources and on the community.

\section{Primary Educational Resources}
Noongar is taught in Western Australian schools. One summary says it is among 24 Aboriginal languages taught at 68 schools to about 10,000 students (2022). Curtin University offers an open online Noongar language and culture course. I saw no URL for either of these, so they are not in the json.

\section{Linguistic and Reference Materials}
\begin{itemize}
\item Wikipedia, ``Noongar language'': overview of dialects, documentation and revitalisation. \url{https://en.wikipedia.org/wiki/Noongar_language}
\item Early and later documentation named in the search summary: George Fletcher Moore's 1842 vocabulary, Wilfrid Douglas's grammar and dictionary, Bindon and Chadwick's \emph{A Nyoongar Wordlist}, and Bernard Rooney's Nyoongar dictionary (Yued area). I saw no URLs for these.
\end{itemize}

\section{Religious and Community Resources}
The Bibbulmun Track site has an Aboriginal interpretation page on Noongar Country and language groups, citing SWALSC, SERCUL and AIATSIS. \url{https://www.bibbulmuntrack.org.au/the-track/aboriginal-interpretation/}
The Noongar Language and Culture Centre (Bunbury, founded 1986) was named in a search summary. I saw no URL for it.

\section{Digital and Audio Resources}
I found no Bibbulman audio or software in search results.

\section{Curriculum Development Context}
Little is specific to Bibbulman. A heritage curriculum would use Noongar materials together with local Bibbulman-country knowledge. It would need to work with Noongar elders and organisations such as SWALSC and the Noongar Language and Culture Centre, and check AIATSIS for archival wordlists.
\end{document}